\begin{abstract}

    \begin{center}
        \large{% Единое место для указания названия работы
Разработка системы шумоподавления и очистки аудиосигналов на основе глубокого обучения
} \\
    \large\textit{Вамбриков Никита Сергеевич} \\[1cm]
    \end{center}

    Работа посвящена исследованию системы восстановления речи на основе
    архитектуры Miipher-2 и самоконтролируемых (SSL) акустических
    представлений. Ключевой компонент оригинальной модели — энкодер USM —
    закрыт и недоступен, поэтому цель работы состоит в подборе и исследовании
    открытых аналогов, обеспечивающих сопоставимое качество очистки речи при
    приемлемых вычислительных затратах.

    Решены следующие задачи: выполнен обзор подходов к улучшению речи;
    воспроизведена схема <<SSL-энкодер — очиститель признаков — вокодер>> с
    открытыми энкодерами (WavLM, HuBERT) и вокодерами (HiFi-GAN, WaveFit);
    исследовано влияние номера слоя извлечения признаков на качество; проведено
    сравнение четырёх комбинаций <<энкодер — вокодер>>.

    Качество оценивалось метриками PESQ, STOI, SI-SDR и DNSMOS, а также по
    вычислительной сложности и скорости. Установлено, что признаки промежуточных
    слоёв WavLM в сочетании с вокодером HiFi-GAN дают наилучший баланс качества
    и быстродействия. Результаты подтверждают применимость открытых
    SSL-фронтендов для построения систем очистки речи и генерации обучающих
    датасетов студийного качества.

    \vfill

\end{abstract}
